\documentclass[11pt,a4paper]{article}
\usepackage[margin=0.85in]{geometry}
\usepackage{amsmath,booktabs,tabularx,microtype}
\usepackage{hyperref}
\hypersetup{colorlinks=true,linkcolor=black,urlcolor=blue,
  pdftitle={CBMR Implementation Comparisons: Expanded Experiment Findings}}
\emergencystretch=3em
\newcolumntype{T}{>{\raggedright\arraybackslash}X}
% Generated by python -m comparison.explanation; do not edit.
\newcommand{\ExpandedRecords}{54}
\newcommand{\ExpandedRuns}{162}
\newcommand{\ExpandedPairs}{18}
\newcommand{\ExpandedFasterOriginal}{4}
\newcommand{\ExpandedFasterAutodiff}{18}
\newcommand{\ExpandedAttempts}{400}
\newcommand{\ExpandedAcceptedFits}{0}
\newcommand{\ExpandedAcceptedInferences}{0}
\newcommand{\ExpandedScoreThreshold}{1.00\times 10^{-4}}
\newcommand{\ExpandedHessianTable}{\begin{tabular}{llrrrr}
\toprule
Size & Model & Original & Improved & Autodiff & Ratio$^\dagger$\\
\midrule
Small & Poisson & 47.20 & 2.06 & 7.36 & 22.93$\times$\\
Small & Aggregated NB & 2.80 & 2.83 & 102.44 & 0.99$\times$\\
Small & Clustered NB & 1.38 & 97.50 & 410.00 & 0.01$\times$\\
Moderate & Poisson & 1.39 & 2.15 & 13.76 & 0.64$\times$\\
Moderate & Aggregated NB & 85.11 & 98.78 & 824.78 & 0.86$\times$\\
Moderate & Clustered NB & 15.67 & 16.90 & 442.64 & 0.93$\times$\\
Large & Poisson & 99.81 & 100.62 & 422.14 & 0.99$\times$\\
Large & Aggregated NB & 4.31 & 3.86 & 508.45 & 1.12$\times$\\
Large & Clustered NB & 1.64 & 2.39 & 196.93 & 0.69$\times$\\
\bottomrule
\end{tabular}}
\newcommand{\ExpandedContrastTable}{\begin{tabular}{llrrrr}
\toprule
Size & Model & Original & Improved & Autodiff & Ratio$^\dagger$\\
\midrule
Small & Poisson & 2.12 & 3.18 & 7.38 & 0.67$\times$\\
Small & Aggregated NB & 3.46 & 4.03 & 115.39 & 0.86$\times$\\
Small & Clustered NB & 2.06 & 2.69 & 9.15 & 0.77$\times$\\
Moderate & Poisson & 1.99 & 96.95 & 347.60 & 0.02$\times$\\
Moderate & Aggregated NB & 102.32 & 98.13 & 719.02 & 1.04$\times$\\
Moderate & Clustered NB & 99.82 & 161.85 & 497.57 & 0.62$\times$\\
Large & Poisson & 2.07 & 3.10 & 107.28 & 0.67$\times$\\
Large & Aggregated NB & 84.32 & 5.01 & 491.25 & 16.81$\times$\\
Large & Clustered NB & 2.31 & 3.10 & 106.81 & 0.75$\times$\\
\bottomrule
\end{tabular}}
\newcommand{\ExpandedCalibrationTable}{\begin{tabular}{lrrrrr}
\toprule
Model & \shortstack{Optimiser\\success} & \shortstack{Accepted\\fit / inference} & \shortstack{Score norm\\minimum} & Median & Maximum\\
\midrule
Poisson & 100 & 0 / 0 & $1.90\times 10^{-3}$ & $6.12\times 10^{-3}$ & $8.43\times 10^{-2}$\\
Independent NB & 100 & 0 / 0 & $2.00\times 10^{-3}$ & $1.08\times 10^{-2}$ & $5.82\times 10^{-2}$\\
Aggregated NB & 100 & 0 / 0 & $2.70\times 10^{-4}$ & $1.07\times 10^{-3}$ & $6.44\times 10^{-3}$\\
Clustered NB & 100 & 0 / 0 & $5.48\times 10^{-4}$ & $2.60\times 10^{-3}$ & $1.97\times 10^{-2}$\\
\bottomrule
\end{tabular}}
\newcommand{\ExpandedMaxAbsolute}{5.87\times 10^{-11}}
\newcommand{\ExpandedMaxRelative}{8.94\times 10^{-15}}
\newcommand{\ExpandedOriginalRSS}{259.5--264.4}
\newcommand{\ExpandedImprovedRSS}{258.7--262.9}
\newcommand{\ExpandedTimingSpread}{2.46, 96.95, 197.17}
\newcommand{\ExpandedBenchmarkDigest}{ef8cf087a589a267}
\newcommand{\ExpandedCalibrationDigest}{77f57e0488d55b1a}

\title{CBMR Implementation Comparisons\\
  \large Motivation, comparison setups and expanded findings}
\author{}
\date{Updated \today}

\begin{document}
\maketitle
\vspace{-1.5em}
\noindent\textbf{Main finding.}
Numerical agreement remains strong and the improved aggregated-NB path is more
stable at extreme parameters. The expanded benchmark does \emph{not} show a
consistent speed advantage over the original, already closed-form implementation.
All \ExpandedAttempts{} expanded calibration attempts were rejected by the
final-score acceptance rule, despite optimiser success flags; consequently
they provide no estimated coverage or false-positive rates.

\medskip
\noindent This updates the earlier smoke-only performance and calibration
summary. Sources are the saved JSON reports in
\texttt{cbmr-proofs/comparison/results/}; experiments were not rerun to produce
this document. The source hashes in both expanded reports match the current
\texttt{cbmr\_code/} and \texttt{cbmr\_improved/} snapshots at generation.

\section{Motivation and comparison setups}
Exact SymPy identities and Lean proofs do not, by themselves, establish
floating-point stability, successful optimisation, faster execution or reliable
statistical inference. The comparison suite addresses these separate questions
without modifying either implementation directory.
GC denotes covariate effects constant across voxels; SV allows those effects
to vary spatially. Covariate values may vary between studies in both designs.

\begin{center}\small
\begin{tabularx}{\textwidth}{>{\raggedright\arraybackslash}p{0.18\textwidth}T}
\toprule
Layer & Setup and purpose\\
\midrule
Likelihoods and derivatives &
Compare objectives, scores and Hessians at identical coefficients and fixed
dispersion. Independent NumPy/SciPy closed formulas and unsimplified float64
PyTorch autodiff provide separate implementations; directional finite
differences provide another check.\\
Models and designs &
Poisson, aggregated NB and clustered NB are tested in both source packages.
Independent-cell NB is reference/helper-only. GC and SV cases include shared
covariates, unequal groups and varying exposures. Fixed quadratic penalties
are external wrappers, not native package features.\\
Covariance &
Compare block/Schur inversion, shared-global cross-group uncertainty, Kronecker
structure and projected sandwiches against dense calculations. Include
positive-definite, indefinite, nearly singular and singular cases.\\
Stability and fits &
Exercise zero counts, extreme means, small dispersion and redundant intercepts.
Identifiable, unpenalised fits use the same starting coefficients and L-BFGS-B
configuration, with dispersion fixed.\\
Performance &
Sequential fresh processes and one numerical thread. Time excludes imports
and setup. Peak RSS is a whole-process high-water mark, not incremental
operation memory.\\
Calibration &
Simulate each model's own distribution and assess coefficient bias, 95\% Wald
coverage, null-moderator false positives and failed attempts separately from
derivative agreement and timing.\\
\bottomrule
\end{tabularx}
\end{center}
\noindent\textbf{Model boundary.}
Aggregated NB is a moment-matched pattern-total likelihood, not independent-cell
NB. Its zero-dispersion limit is the \emph{aggregate} Poisson objective.
Global curvature need not equal full cell-level Poisson because study-allocation
information is omitted.

\clearpage
\section{Established numerical agreement and stability}
\subsection{Regression, covariance and complete fits}
The saved \texttt{regression.json} records \textbf{512 required comparisons
with zero failures}, plus 47 diagnostic observations. The standalone covariance
and fit reports also have zero required failures (46 and 129 comparisons);
these overlap the full regression report and must not be added to it.

\begin{itemize}
\item Across 312 derivative comparisons, the largest recorded absolute
discrepancy was approximately $5.1\times10^{-11}$ and the largest normalised
Frobenius discrepancy was $3.1\times10^{-12}$.
\item Covariance algebra agreed closely with dense calculations, with maximum
absolute discrepancy approximately $1.4\times10^{-14}$. A nearly singular
example had condition number approximately $6.9\times10^{12}$: small backward
error does not guarantee accurate inverse entries. A redundant-intercept
example exposed five identifiable dimensions for six coefficients.
\item All 28 small implementation/model/design fits reported optimiser success;
all 100 fit-agreement comparisons passed. The largest final score infinity
norm was $4.23\times10^{-6}$, below the required $2\times10^{-5}$.
\end{itemize}
Most derivative comparisons require elementwise agreement with
\texttt{rtol=1e-8}, \texttt{atol=1e-10}; finite-difference and fit checks use their
own stated thresholds. These are finite numerical checks, not additional
Lean proofs or statistical-calibration results.

\subsection{Aggregated-NB stability}
At spatial log means $-1000$ and $+1000$, the improved implementation returned
finite objectives, scores and Hessians, whereas the original implementation
and unsimplified reference produced nonfinite results.
At dispersion $10^{-12}$, the improved Hessian approached the correct
aggregate-Poisson limit with maximum absolute error $9.45\times10^{-11}$,
versus approximately $0.097$ for the original and $1.34$ for the reference.
Autodiff is therefore not treated as an infallible floating-point oracle.
Other difficult regimes are retained as diagnostics, rather than hidden behind
the zero-failure headline.

\section{Expanded performance experiment}
\subsection{What changed from the smoke benchmark}
\texttt{benchmark-expanded.json} adds all three source likelihoods, two workloads
and three problem sizes, using GC designs and three fresh-process repetitions
per configuration. All \ExpandedRecords{} implementation/configuration records
completed successfully: \ExpandedRuns{} measured workers plus 18 separate
reference workers. The 18 size/model/workload cases are each run with the
original, improved and unsimplified-autodiff implementations.

\begin{center}\small
\begin{tabular}{lrrrrr}
\toprule
Size & Experiments & Voxels & Bases & Groups & Parameters\\
\midrule
Small & 32 & 32 & 4 & 3 & 13\\
Moderate & 64 & 96 & 8 & 4 & 33\\
Large & 128 & 256 & 8 & 4 & 33\\
\bottomrule
\end{tabular}
\end{center}
Small-to-moderate scaling changes both data size and parameter count.
Moderate-to-large keeps 33 parameters fixed. Even ``large'' is a bounded
synthetic problem, not a whole-brain production benchmark.

\clearpage
\subsection{Measured runtimes}
All entries below are \textbf{median milliseconds across three repetitions}.
The source implementations already contain closed-form Hessians; this tests
the additional implementation changes, not merely closed forms versus autodiff.
The autodiff benchmark uses a row-wise reverse-mode Hessian, not a reproduction
of production forward-mode memory behaviour.

\paragraph{Hessian assembly.}
\begin{center}\small
\ExpandedHessianTable
\end{center}

\paragraph{Contrast covariance.}
The same two-row contrast selects a group-intercept difference and a global
moderator. The original route assembles information and a full covariance;
the improved route assembles structured information and uses projected solves.
\begin{center}\small
\ExpandedContrastTable
\end{center}
\noindent$^\dagger$\textbf{Ratio = original median / improved median.}
Above one means the improved route was faster; below one means it was slower.
Ratios are descriptive measurements, not statistically established speedups.

\subsection{Interpretation and limitations}
The improved implementation was faster than the original in
\textbf{\ExpandedFasterOriginal{} of \ExpandedPairs{} paired cases},
and faster than the unsimplified-autodiff route in
\ExpandedFasterAutodiff{} of \ExpandedPairs{} cases.
There is \textbf{no consistent additional speed advantage over the original}
in this bounded experiment.

Numerical agreement held throughout: across original and improved benchmark
outputs, the largest absolute discrepancy was $\ExpandedMaxAbsolute$ and the
largest normalised Frobenius discrepancy was $\ExpandedMaxRelative$.
This supports implementation agreement in these cases, independently of
which method happened to be faster.

The median whole-process RSS ranges were \ExpandedOriginalRSS{} MiB for the
original and \ExpandedImprovedRSS{} MiB for the improved implementation.
Interpreter, imports and setup dominate these totals; the results do not
demonstrate meaningful operation-level or production-scale memory savings.

Only three repetitions were used, without warmup. For example, the improved
moderate Poisson contrast workload took \ExpandedTimingSpread{} ms in its
three runs. This substantial variation limits timing conclusions and
precludes a reliable scaling claim from these measurements alone.

\clearpage
\section{Expanded calibration: no accepted inference replicates}
\texttt{calibration-expanded.json} attempted \textbf{100 replicates per model}
on the large GC design (128 experiments, 256 voxels, eight bases, four groups,
33 coefficients), for \ExpandedAttempts{} attempts in total.
Dispersion stayed fixed at its true value; all fits started from zero
coefficients with the same L-BFGS-B protocol.
Each model was simulated from its own specified sampling distribution.

\begin{center}\small
\ExpandedCalibrationTable
\end{center}
\noindent Counts are out of 100 attempts per row. Score norm means
$\|\nabla(-\ell)\|_\infty$. The acceptance rule requires a successful optimiser
flag, finite results and score norm at most $\ExpandedScoreThreshold$ before
attempting inference.

\subsection{Why the report says ``completed with failures''}
All \ExpandedAttempts{} optimisations returned a success flag and terminated
with the relative-objective-reduction message
\texttt{REL\_REDUCTION\_OF\_F\_<=\_FACTR*EPSMCH}.
Nevertheless, \textbf{every final score norm exceeded the separate
$\ExpandedScoreThreshold$ requirement}. Thus there were
\ExpandedAcceptedFits{} accepted fits and \ExpandedAcceptedInferences{}
accepted inference replicates.
The recorded failure is an \emph{acceptance-criterion failure}, not evidence
that all optimisers crashed or that all information matrices were singular.
Inference was not attempted after those rejections.

\textbf{Coverage, false-positive rates and the report's accepted-fit bias
summaries are unavailable, not zero.}
Their denominators are zero and the corresponding JSON values are
\texttt{null}. The run therefore cannot establish either good or poor
statistical calibration.
The earlier three-replicate-per-model smoke experiment had 12 accepted fits,
but was far too small to support calibration claims; it does not override
the expanded run's acceptance failures.

\subsection{What can be concluded and what remains to do}
The expanded experiment strengthens the evidence for numerical agreement,
but does not demonstrate a universal performance gain or validate inferential
calibration. Before using the calibration experiment substantively, investigate
why objective-reduction stopping precedes the pre-specified score requirement,
validate an appropriate fitting/stopping protocol and rerun the simulations.
Do not silently relax the rule, count rejected fits as successful, or
interpret missing calibration estimates as zero error rates.

\section{Provenance and reproducibility}
Individual timings, optimiser diagnostics and rejected attempts remain in the
saved JSON reports under \texttt{comparison/results/}. Report fingerprints
(first 16 SHA-256 digits) are:
\begin{center}\small
\begin{tabular}{ll}
\toprule
Report & Fingerprint\\
\midrule
\texttt{benchmark-expanded.json} & \texttt{\ExpandedBenchmarkDigest}\\
\texttt{calibration-expanded.json} & \texttt{\ExpandedCalibrationDigest}\\
\bottomrule
\end{tabular}
\end{center}
The editable source is \texttt{docs/explanation.tex}.
\texttt{make explanation} regenerates its expanded-result tables from the saved
JSON reports and rebuilds \texttt{explanation.pdf}; it does not rerun experiments.
The generator checks source snapshots and the assumptions behind this
interpretation, stopping for review if the experiment outcomes change.

\end{document}
